\documentclass[12pt]{article}
\usepackage[utf8]{inputenc}
\usepackage[T1]{fontenc}
\usepackage[margin=1in]{geometry}
\usepackage{setspace}
\usepackage{xcolor}
\usepackage[colorlinks=true, urlcolor=blue!60!black, linkcolor=blue!60!black]{hyperref}
\setlength{\parindent}{0em}
\setlength{\parskip}{0.6em}
\renewcommand{\baselinestretch}{1.0}
\begin{document}
\begin{center}
{\LARGE \textit{Trump v. Barbara}, 609 U.S. \_\_\_ (2026) (No. 25-365)\par}
\vspace{0.5em}
{\large Case Brief --- Two-Page Edition\par}
\end{center}

\section{Memory Jogger}
The Citizenship Clause guarantees birthright citizenship to U.S.-born children of unlawfully or temporarily present parents; the President's contrary order fell (\href{https://www.law.cornell.edu/supremecourt/text/25-365}{Cornell LII}).

\section{Detailed Case Facts}
Executive Order No.~14160 denied citizenship to U.S.-born children whose mothers were unlawfully or temporarily present and whose fathers were not citizens or permanent residents, claiming they were not ``subject to the jurisdiction'' of the United States. After \textit{Trump v.\ CASA} limited universal injunctions, the ACLU filed a class action led by a Honduran mother, ``Barbara.''

\section{Procedural History}
The District of New Hampshire certified a nationwide class of affected children and enjoined the Order; the Supreme Court granted certiorari before judgment and affirmed on June 30, 2026.

\section{Judicial Votes}
Affirmed, 6--3 (5--4 on the Constitution). Roberts wrote for Sotomayor, Kagan, Barrett, and Jackson; Jackson concurred; Kavanaugh concurred in the judgment on statutory grounds; Thomas (with Gorsuch), Alito, and Gorsuch dissented.

\section{Holding}
``Subject to the jurisdiction'' means subject to the United States' territorial power to govern, incorporating common-law \textit{jus soli} with narrow exceptions such as diplomats' children. Neither the President nor Congress may deny citizenship the Clause guarantees.

\section{Analysis of Opinions}
\textbf{Majority:} \textit{Calvin's Case}, the repudiation of \textit{Dred Scott}, \textit{The Schooner Exchange}, and \textit{Wong Kim Ark} confirm the territorial rule; domicile theories are ``dramatically revisionist.'' \textbf{Jackson:} the Clause was deliberately universal. \textbf{Kavanaugh:} \S\ 1401(a) bars the Order, but Congress could add exceptions. \textbf{Thomas:} citizenship requires domicile. \textbf{Alito:} the ruling rewards ``birth tourists.''

\section{Examples: Future Applications}
\textbf{Same-Side:} a statute denying citizenship to children of parents who entered without inspection is unconstitutional. \textbf{Opposite-Side:} a child of accredited foreign ambassadors is not a citizen. \textbf{Fence-Sitter:} a child born in American Samoa, where the \textit{Insular Cases} still govern whether it is ``in the United States.''

\section{Critique}
Originalists split: Barnett and Ramsey backed the majority, while critics drawing on Lash and Wurman said it read ``jurisdiction'' out of the Clause. Frost and Rosenbloom warned that the guarantee held by one vote.

\section{Key Quotations}
\begin{quote}
``The Citizenship Clause uses jurisdiction in its ordinary sense---referring to the power of the United States to govern those within its territory.''
\end{quote}
\noindent --- Chief Justice Roberts, Opinion of the Court.

\end{document}
